\chapter[Neural State-Space Model Applications]
{Neural Networks and State-Space Model Applications}
\label{ch:bf-neural-network-ssm-applications}

Neural state-space models combine a learned nonlinear transition with an
explicit probabilistic state and observation model.  This combination is
useful when the transition must retain memory that is richer than a low-order
linear recursion, but the application still requires filtering, forecasting,
likelihood-based parameter inference, and uncertainty propagation.  It also
creates a validation problem.  A neural transition can be strongly
overparameterized, so two posterior approximations may differ in parameter
space while inducing nearly the same forecast distribution.  Conversely, two
samplers can agree on a few parameter summaries and still disagree on the
predictive quantities for which the model is used.

This chapter develops a running BayesFilter application based on a stochastic
state-space LSTM (SSL-LSTM).  It then defines a posterior-predictive
equivalence criterion for comparing ordinary HMC with exact-corrected
neural-transport HMC.  The model description follows the current BayesFilter
TensorFlow implementation.  The predictive-equivalence criterion is a
BayesFilter \code{extension\_or\_invention}; it is not a Zhao--Cui
source-faithfulness claim.

\section{A Stochastic State-Space LSTM}
\label{sec:bf-ssl-lstm-model}

Let $z_t\in\R^k$ denote the stochastic latent state, $a_t\in\R^h$ the LSTM
hidden state, and $c_t\in\R^h$ the LSTM cell state.  The augmented filtering
state is
\begin{equation}
  x_t = (z_t^\top,a_t^\top,c_t^\top)^\top
  \in \R^{k+2h}.
  \label{eq:bf-ssl-lstm-state}
\end{equation}
Only $z_t$ receives a new structural innovation.  Conditional on the previous
state, $a_t$ and $c_t$ are deterministic completions.  This partition is
important for sigma-point filtering because it avoids pretending that the
memory coordinates contain independent process noise.

For gate $g\in\{i,f,o,r\}$, define
\begin{align}
  \eta_{g,t}
    &= W_g z_{t-1}+U_g a_{t-1}+b_g,                                      \\
  i_t &= \operatorname{sigmoid}(\eta_{i,t}),                              \\
  f_t &= \operatorname{sigmoid}(\eta_{f,t}),                              \\
  o_t &= \operatorname{sigmoid}(\eta_{o,t}),                              \\
  r_t &= \tanh(\eta_{r,t}).
  \label{eq:bf-ssl-lstm-gates}
\end{align}
Here
\[
  W_g\in\R^{h\times k},\qquad
  U_g\in\R^{h\times h},\qquad
  b_g\in\R^h,
\]
so every gate vector lies in $\R^h$.  The latent and observation maps satisfy
$A\in\R^{k\times h}$, $d\in\R^k$, $C\in\R^{d_y\times k}$, and
$e\in\R^{d_y}$.
The candidate gate is denoted by $r$ to avoid confusing it with the cell state.
The memory recursion is
\begin{align}
  c_t &= f_t\odot c_{t-1}+i_t\odot r_t,                                    \\
  a_t &= o_t\odot\tanh(c_t).
  \label{eq:bf-ssl-lstm-memory}
\end{align}
The stochastic latent transition and observation law are
\begin{align}
  z_t &= A a_t+d+\Sigma_z^{1/2}\epsilon_t,
  &\epsilon_t&\sim\calN(0,I_k),                                           \\
  y_t &= C z_t+e+\Sigma_y^{1/2}\nu_t,
  &\nu_t&\sim\calN(0,I_{d_y}),
  \label{eq:bf-ssl-lstm-state-observation}
\end{align}
with independent innovations across time and between the two equations.  The
current implementation uses diagonal $\Sigma_z$ and $\Sigma_y$.  Initial
states follow
\begin{equation}
  x_0\sim\calN(m_0,\diag(s_0^2)).
  \label{eq:bf-ssl-lstm-initial-law}
\end{equation}
Positive standard deviations are constructed from unconstrained coordinates
by a softplus transform plus a small positive floor.

Equations~\eqref{eq:bf-ssl-lstm-gates}--\eqref{eq:bf-ssl-lstm-state-observation}
have the usual LSTM memory mechanism \citep{hochreiter1997lstm}, but they are
embedded in a generative state-space model rather than used only as a
deterministic predictor.  The transition produces a probability law, the
filter integrates over latent uncertainty, and Bayesian inference propagates
parameter uncertainty into future observations.

\subsection{Parameter chart}
\label{subsec:bf-ssl-lstm-parameter-chart}

The unconstrained BayesFilter parameter vector contains
\begin{equation}
  \theta =
  \left(
    \{W_g,U_g,b_g\}_{g\in\{i,f,o,r\}},
    A,d,C,e,m_0,\widetilde s_0,
    \widetilde\sigma_z,\widetilde\sigma_y
  \right).
  \label{eq:bf-ssl-lstm-parameter-chart}
\end{equation}
For dimensions $(k,h,d_y)$, the implementation stores the four input-gate
matrices, four recurrent matrices, four biases, the latent and observation
maps, the initial mean and standard deviations, and the process and
observation standard deviations in one fixed order.  Static dimensions and a
fixed chart are part of the compiled filtering and HMC target contract.

\subsection{Structural SVD-UKF likelihood}
\label{subsec:bf-ssl-lstm-svd-ukf-likelihood}

The running application evaluates the likelihood with BayesFilter's structural
SVD sigma-point filter.  Sigma-point integration is performed in the
$k$-dimensional innovation space.  Each innovation point is lifted through
Equations~\eqref{eq:bf-ssl-lstm-gates}--\eqref{eq:bf-ssl-lstm-state-observation},
and the LSTM hidden and cell states are completed deterministically.  This
preserves the state partition in Equation~\eqref{eq:bf-ssl-lstm-state}.

Let $\ell_{\mathrm{SVD\text{-}UKF}}(\theta;y_{1:T})$ denote the resulting
filtering log likelihood.  The posterior target is
\begin{equation}
  \pi(\theta\mid y_{1:T})
  \propto
  \exp\!\left\{\ell_{\mathrm{SVD\text{-}UKF}}
    (\theta;y_{1:T})\right\}p(\theta).
  \label{eq:bf-ssl-lstm-filtering-posterior}
\end{equation}
This is an exact scalar target relative to the fixed SVD-UKF recursion, but the
SVD-UKF likelihood is an approximation to the nonlinear model likelihood.
Sampler validation must not erase that distinction.

\section{The Scalar Four-Parameter Application}
\label{sec:bf-scalar-ssl-lstm-application}

The scalar validation problem sets
\begin{equation}
  k=h=d_y=1,
  \qquad \dim(x_t)=3,
  \qquad T=30.
  \label{eq:bf-scalar-ssl-lstm-dimensions}
\end{equation}
The full scalar chart has 24 entries.  The current validation target holds 20
entries fixed at the inherited fixture values and estimates
\begin{equation}
  \psi=(A,d,C,e)^\top\in\R^4.
  \label{eq:bf-scalar-ssl-lstm-free-parameters}
\end{equation}
In repository names these are
\begin{itemize}
  \item \code{latent\_mean\_weight.0.0};
  \item \code{latent\_mean\_bias.0};
  \item \code{observation\_weight.0.0};
  \item \code{observation\_bias.0}.
\end{itemize}
The prior used by the scalar filtering target is
\begin{equation}
  \psi\sim\calN(\psi_0,4^2 I_4),
  \label{eq:bf-scalar-ssl-lstm-prior}
\end{equation}
where $\psi_0$ is the fixed simulation center.  The observation path is
generated with a fixed stateless seed and is shared by every sampler arm.

This restriction creates a deliberately small but non-Gaussian filtering
target.  It is useful for testing geometry, transport, HMC, and predictive
validation without first solving full 24-parameter identification.  It does
not establish that the four-parameter restriction is scientifically adequate,
that the remaining neural-network parameters are known, or that the fixture is
source-faithful to an external SSL-LSTM implementation.

\section{Why Parameter Agreement Is Not the Only Estimand}
\label{sec:bf-predictive-vs-parameter-agreement}

Traditional sampler checks often compare posterior parameter means,
variances, quantiles, or an independently computed parameter posterior.  Those
checks remain useful when the parameter posterior itself is the scientific
object.  They can be unnecessarily strong for prediction when multiple
parameter configurations induce nearly the same transition and observation
law.  Neural networks make this issue especially visible because symmetries,
weak identification, and compensating weights can create broad or curved
parameter sets with similar functional behavior.

The converse warning is equally important.  Predictive agreement is weaker
than equality of parameter posteriors.  If $F$ maps a parameter and filtered
state distribution into a forecast law, then
\begin{equation}
  F_\#\pi_A = F_\#\pi_B
  \quad\not\Longrightarrow\quad
  \pi_A=\pi_B.
  \label{eq:bf-predictive-not-parameter-identification}
\end{equation}
The pushforward can discard directions that matter for structural
interpretation even when they do not matter for the selected forecast.  The
appropriate claim is therefore \emph{predictive-functional equivalence}, not
full posterior correctness or parameter identification.

\section{Posterior-Predictive Law}
\label{sec:bf-ssl-lstm-posterior-predictive-law}

Fix the observed history $y_{1:T}$ and a maximum forecast horizon $H$.  For
sampler arm $s$, let $\pi_s(d\psi,dx_T\mid y_{1:T})$ be the joint parameter and
terminal-state approximation induced by the fixed filtering target and that
arm's posterior draws.  Its joint predictive path law is
\begin{equation}
  P_s^H(d\widetilde y)
  = \int
      p_\psi(d\widetilde y_{T+1:T+H}\mid x_T)
      \pi_s(d\psi,dx_T\mid y_{1:T}).
  \label{eq:bf-ssl-lstm-joint-predictive-law}
\end{equation}
Every sampler comparison must use the same data, parameter chart, terminal
filter convention, forecast equations, horizon, and innovation law.  A frozen
stateless innovation bank may be shared between arms to reduce conditional
simulation noise, but uncertainty calculations must retain the chain and
forecast-cluster dependence that this design creates.  Because this coupling
is not part of either marginal predictive law, the omnibus distributional gate
should use independently generated arm-specific innovation banks; the shared
bank is a paired precision and robustness diagnostic.

\subsection{Horizon-specific moments}
\label{subsec:bf-ssl-lstm-predictive-moments}

For $h=1,\ldots,H$, define
\begin{align}
  \mu_{s,h}
    &= \E_s[Y_{T+h}\mid y_{1:T}],                                         \\
  v_{s,h}
    &= \Var_s(Y_{T+h}\mid y_{1:T}),                                      \\
  \gamma_{s,h}^{(r)}
    &= \E_s[(Y_{T+h}-\mu_{s,h})^r\mid y_{1:T}],\qquad r\geq 3,            \\
  \Gamma_{s,hk}
    &= \Cov_s(Y_{T+h},Y_{T+k}\mid y_{1:T}).
  \label{eq:bf-ssl-lstm-predictive-moments}
\end{align}
The first central moment is identically zero, so the first comparison is the
raw predictive mean.  Variance is the second central moment.  Third and fourth
central moments can reveal skewness and tail-shape differences, but their Monte
Carlo uncertainty is substantially larger.  Cross-horizon covariance matters
because equality of all marginal forecast laws does not imply equality of the
joint path law.

When the filter can propagate conditional means and covariances
deterministically for each posterior draw, the law of total expectation and
variance gives a lower-noise estimator:
\begin{align}
  \mu_{s,h}
    &= \E_{\pi_s}[m_h(\psi,x_T)],                                         \\
  v_{s,h}
    &= \E_{\pi_s}\!\left[V_h(\psi,x_T)+m_h(\psi,x_T)^2\right]
       -\mu_{s,h}^2.
  \label{eq:bf-ssl-lstm-rao-blackwell-moments}
\end{align}
BayesFilter should prefer these Rao--Blackwellized moments when the forecast
recursion and state approximation make them available.  Simulated forecast
paths are still required for general path-distribution comparisons.

\section[MGFs and Stable Distribution Features]
{From Moment Generating Functions to Stable Distribution Features}
\label{sec:bf-ssl-lstm-mgf-cf}

Whenever the MGF is finite on an open neighborhood of zero and termwise
differentiation is justified, its Taylor expansion is
\begin{equation}
  M_s(t)=\E_s[\exp(tY)]
  =\sum_{r=0}^{\infty}\frac{t^r}{r!}\E_s[Y^r]
  \label{eq:bf-ssl-lstm-mgf}
\end{equation}
This expansion provides an attractive formal connection between a weighted
collection of moments and the full predictive distribution.  It is not,
however, a robust default empirical weight.  The MGF need not be finite away
from zero, and a small number of tail draws can dominate $\exp(tY)$.

The characteristic function is safer:
\begin{equation}
  \varphi_s(\omega)
  =\E_s[\exp(i\omega^\top\widetilde Y)]
  =\E_s[\cos(\omega^\top\widetilde Y)]
   +i\E_s[\sin(\omega^\top\widetilde Y)].
  \label{eq:bf-ssl-lstm-characteristic-function}
\end{equation}
It always exists, its empirical features are bounded, and the joint
characteristic function identifies the joint predictive law.  A weighted
distance is
\begin{equation}
  D_w(P_A^H,P_B^H)
  =\int
    \left|\varphi_A(\omega)-\varphi_B(\omega)\right|^2
    w(d\omega).
  \label{eq:bf-ssl-lstm-weighted-cf-distance}
\end{equation}
For suitable spectral weights, this is the same population discrepancy as a
kernel maximum mean discrepancy (MMD) \citep{gretton2012kernel,
sriperumbudur2010hilbert}.  A Gaussian spectral weight corresponds to a
Gaussian radial-basis kernel.  A fixed mixture of bandwidths is preferable to
choosing a bandwidth after inspecting which choice makes two samplers agree.

\section{What the Weights Mean}
\label{sec:bf-ssl-lstm-predictive-weights}

Three different weights enter a predictive comparison and should not be
collapsed into one tuning parameter.

\paragraph{Scientific horizon weights.}
Weights $\lambda_h\geq0$ with $\sum_h\lambda_h=1$ state whether near-term and
long-term forecasts have equal scientific value.  Equal horizon weights are
the neutral default.  Discounted weights are legitimate only when the
application declares that shorter horizons matter more; they must not be
estimated to hide a long-horizon discrepancy.

\paragraph{Feature scaling.}
Means, variances, fourth moments, and cross-products have different units.
They must be standardized by a frozen external scale or by a baseline-only
pilot whose draws are excluded from confirmation.  Scaling is a numerical and
interpretive operation, not evidence that a discrepancy is small.

\paragraph{Statistical precision weights.}
Let $\widehat\Delta$ collect real and imaginary empirical characteristic
features, or a fixed vector of standardized moment differences.  If
$\widehat\Delta\in\R^q$, $\widehat S\in\R^{q\times q}$ is a symmetric
positive-semidefinite estimate of its long-run covariance, and $\lambda>0$,
then $\widehat S+\lambda I_q$ is positive definite.  The regularized quadratic
form
\begin{equation}
  Q_\lambda
  =\widehat\Delta^\top
    (\widehat S+\lambda I_q)^{-1}\widehat\Delta
  \label{eq:bf-ssl-lstm-gmm-weight}
\end{equation}
uses the efficient GMM weighting principle for a fixed feature vector
\citep{hansen1982gmm}.  Here $\widehat S$ must account for MCMC
autocorrelation, separate chains, repeated forecast innovations, and any
common-random-number coupling.  The ridge $\lambda$ and feature grid must be
chosen in calibration and frozen before the confirmatory arm comparison.

Inverse-covariance weighting does not replace scientific equivalence margins.
A very precisely estimated but scientifically negligible feature should not
dominate the decision merely because its Monte Carlo variance is small.

\section{An Equivalence Gate, Not a Failure-to-Reject Gate}
\label{sec:bf-ssl-lstm-predictive-equivalence-gate}

Let $\delta_{\mu,h}$ be a predictive-scale-standardized mean difference and
let $\delta_{\log v,h}$ be the difference in log predictive variance.  The
log scale makes the variance comparison dimensionless and treats reciprocal
variance ratios symmetrically.  Let $\mathcal I_{\mu,h}$ and
$\mathcal I_{\log v,h}$ denote their simultaneous confidence intervals.  The
primary interpretable gate requires
\begin{align}
  \mathcal I_{\mu,h}
    &\subset (-\varepsilon_{\mu,h},\varepsilon_{\mu,h}),                  \\
  \mathcal I_{\log v,h}
    &\subset (-\varepsilon_{\log v,h},\varepsilon_{\log v,h}),
  \qquad h=1,\ldots,H.
  \label{eq:bf-ssl-lstm-component-equivalence}
\end{align}
The omnibus gate additionally requires an uncertainty upper bound for
$D_w$ or MMD to lie below a predeclared practical tolerance.  Merely failing to
reject equality is insufficient: short chains and noisy high-order moments can
make a materially different pair of predictive laws look statistically
indistinguishable.

Mean and log variance are the initial co-primary interpretable features.
Third and fourth central moments, quantiles, and cross-horizon covariance begin
as explanatory diagnostics.  They can become promotion criteria only after a
reviewed calibration shows adequate finite-sample power and stable uncertainty
estimation.

\section{Ordinary HMC and Exact-Corrected NeuTra-HMC}
\label{sec:bf-ssl-lstm-neutra-comparison}

Let $u$ denote the MAP-local coordinates of the four-parameter filtering
target, and let $T_\phi:\R^4\rightarrow\R^4$ be an invertible learned
transport.  NeuTra trains $T_\phi$ so that a standard Gaussian base is pushed
toward the target geometry \citep{hoffman2019neutra}.  After training, the
transport is frozen and HMC targets
\begin{equation}
  \log\pi_z(z)
  =\log\pi_u(T_\phi(z))
   +\log\left|\det J_{T_\phi}(z)\right|.
  \label{eq:bf-ssl-lstm-exact-neutra-target}
\end{equation}
The Jacobian correction and Metropolis correction preserve the intended
filtering target even when the learned flow is imperfect, subject to the usual
numerical and implementation checks.  Training loss, validation loss, or raw
flow samples are not posterior-correctness evidence.

The functional comparison is therefore
\begin{equation}
  \text{ordinary MAP-local HMC}
  \quad\text{versus}\quad
  \text{frozen-transport exact-corrected NeuTra-HMC},
  \label{eq:bf-ssl-lstm-sampler-comparison}
\end{equation}
with both arms passed through the same forecast operator in
Equation~\eqref{eq:bf-ssl-lstm-joint-predictive-law}.  Each arm must be tuned
in its own geometry.  Reusing the ordinary-HMC step size or trajectory length
for NeuTra-HMC is not a fair comparison.

\section{Three Evidence Ledgers}
\label{sec:bf-ssl-lstm-three-ledgers}

\begin{center}
\begin{tabular}{p{0.23\textwidth}p{0.31\textwidth}p{0.34\textwidth}}
\toprule
Ledger & Question & Required evidence \\
\midrule
Engineering correctness & Does the implementation evaluate the intended
model and forecast? & Equation-to-code checks, deterministic recursions,
linear-Gaussian forecast oracles, scalar/batch parity, transform round trips,
finite values, and XLA parity. \\
Sampler validity & Does each chain explore its declared target? & Native
divergence telemetry, finite log probabilities and gradients, acceptance,
movement, chain-aware R-hat/ESS/MCSE, independent seeds, and frozen artifacts.
\\
Predictive-functional equivalence & Do the two valid samplers induce the same
practically relevant forecast law? & Simultaneous mean/log-variance
equivalence, joint-path characteristic-function or MMD equivalence, and
independent confirmation with frozen weights and tolerances. \\
\bottomrule
\end{tabular}
\end{center}

A failure in the first or second ledger invalidates interpretation of the
third.  A pass in the third ledger does not retroactively prove the first two.
Two samplers can agree because they share a target bug or a forecast bug.
Independent forecast-oracle tests and, where feasible, repeated synthetic-data
calibration remain necessary.

\section{Model Adequacy Is a Separate Question}
\label{sec:bf-ssl-lstm-model-adequacy}

Predictive equivalence asks whether two posterior computations give the same
forecast law.  Predictive accuracy asks whether that law describes future
observations.  Held-out log predictive density, CRPS, coverage, and calibration
are proper model-evaluation tools \citep{gneiting2007proper}, but they do not
replace the sampler-equivalence gate.  The two questions should use separate
artifacts:
\begin{enumerate}
  \item same-data, same-model predictive equivalence for computation;
  \item held-out or rolling-origin proper scores for model adequacy;
  \item parameter or structural checks when scientific interpretation depends
    on identification rather than prediction alone.
\end{enumerate}

\section{Claim Boundary}
\label{sec:bf-ssl-lstm-predictive-claim-boundary}

A successful predictive-equivalence result may support the statement that two
validated sampler arms induce practically equivalent $H$-step predictive laws
for the fixed scalar SSL-LSTM filtering target, data set, forecast operator,
and declared feature/tolerance design.  It does not by itself establish:
\begin{itemize}
  \item equality or correctness of the parameter posterior;
  \item global parameter identification;
  \item exactness of the SVD-UKF approximation to the nonlinear likelihood;
  \item empirical adequacy on unseen data;
  \item broad HMC, NeuTra, GPU, or default-method readiness;
  \item sampler superiority or a statistically supported ranking;
  \item Zhao--Cui source faithfulness.
\end{itemize}

The implementation and confirmatory test program for this chapter is specified
in the predictive-equivalence master plan under \code{docs/plans}.  Until its
forecast, sampler, calibration, and independent-confirmation gates pass, this
chapter defines a research design rather than reporting a completed empirical
result.
